\chapter{统计推断与运筹优化理论}\label{ch:theory}

\section{推断目标分化：瞬态参数与跨期平稳基准}\label{sec:targets}

在第二章建立的两阶段分层生成模型下，本章进入本文的核心理论推导。以往传染病监测理论之所以出现“盲目预设规模越大越好”的困境，其数理根源在于混淆了两个截然不同的统计推断目标：

\begin{definition}[\textbf{目标 A：当期瞬时再生数 $R_t$ 的点估计}]\label{def:target_A}
在特定时间代际 $t$（如某一特定自然周），全社会暴露于当期特定的气象条件、人群聚集水平与临时管控政策之下。此时，瞬态再生数 $R_t$ 是一个已实现的确定性瞬时状态，统计推断的目标是利用当期观测报告病例数 $C_t$ 精确测量该瞬态强度 $R_t$。
\end{definition}

\begin{definition}[\textbf{目标 B：跨期平稳基础传播基准 $\bar{R}$ 的估计与外推}]\label{def:target_B}
公共卫生中长期资源储备、疫苗采购与多期流行趋势研判，需要评估的是消除单周偶发扰动后的基础传播趋势基准 $\bar{R} \triangleq \E[R_t \mid n]$。此时，全社会共享的宏观环境变迁表现为不可预测的外生随机扰动，满足条件方差 $\Var(R_t \mid n) = v_R > 0$。推断目标是利用有限的监测数据估计该潜在基线 $\bar{R}$。
\end{definition}

本章将严格证明：推断目标由 A 切换至 B，将引发方差标度律的根本性断裂！

\section{目标 A 的点估计与 Cramér--Rao 渐近有效性}\label{sec:prop1}

本节针对推断目标 A，推导其极大似然估计量，并证明其无偏性与达到 Cramér--Rao 下界的优良性质。

\begin{proposition}[\textbf{命题 1：条件似然、无偏性与 UMVUE 达到}]\label{prop:prop1}
设已知上一代风险集真实规模 $n$、固定报告率 $\rho$ 及离散参数 $k$。
\begin{enumerate}[label=(\roman*)]
    \item 观测样本 $C_t$ 对瞬态参数 $R_t$ 的极大似然估计量（MLE）与矩估计量严格重合，由下式给出：
    \begin{equation}
    \widehat{R}_t^{\text{MLE}} = \frac{C_t}{\rho n} = \frac{C_t}{m}, \label{eq:estimator_def}
    \end{equation}
    其中 $m \triangleq \rho n$ 为期望报告规模；
    \item 估计量 $\widehat{R}_t$ 在全域 $R_t \in (0, \infty)$ 上严格无偏：$\E[\widehat{R}_t \mid n, R_t] = R_t$；
    \item 在内部参数空间中，估计量 $\widehat{R}_t$ 严格达到 Cramér--Rao 下界，为一致最小方差无偏估计量（UMVUE），其条件方差为：
    \begin{equation}
    \Var(\widehat{R}_t \mid n, R_t) = \mathcal{I}(R_t \mid n)^{-1} = \frac{R_t}{\rho n} + \frac{R_t^2}{k n} = \mathcal{O}\left( \frac{1}{m} \right), \label{eq:var_prop1}
    \end{equation}
    其中 $\mathcal{I}(R_t \mid n)$ 为关于参数 $R_t$ 的 Fisher 信息量。
\end{enumerate}
\end{proposition}

\begin{proof}
下面给出保姆级分步证明：

\textbf{第一步：建立对数似然方程与求解得分函数（Score Function）}

由定理~\ref{thm:closure}，观测样本 $C_t$ 服从负二项分布 $\NB(nk, p^*)$，其中 $p^* = \frac{k}{k + \rho R_t}$，$1 - p^* = \frac{\rho R_t}{k + \rho R_t}$。记观测值为 $c \in \{0, 1, 2, \dots\}$，似然函数为：
\begin{equation}
L(R_t; c) = \frac{\Gamma(c + nk)}{c! \Gamma(nk)} (p^*)^{nk} (1 - p^*)^c = \frac{\Gamma(c + nk)}{c! \Gamma(nk)} \left( \frac{k}{k + \rho R_t} \right)^{nk} \left( \frac{\rho R_t}{k + \rho R_t} \right)^c.
\end{equation}
对似然函数取自然对数 $\ln$：
\begin{align}
\ln L(R_t) &= \ln \left[ \frac{\Gamma(c + nk)}{c! \Gamma(nk)} \right] + nk \ln k - nk \ln(k + \rho R_t) + c \ln \rho + c \ln R_t - c \ln(k + \rho R_t) \notag \\
&= \underbrace{\ln \left[ \frac{\Gamma(c + nk)}{c! \Gamma(nk)} \right] + nk \ln k + c \ln \rho}_{\text{与未知参数 } R_t \text{ 完全无关的常数，简记为 } \text{const}} + c \ln R_t - (nk + c) \ln(k + \rho R_t). \label{eq:log_likelihood}
\end{align}

对参数 $R_t$ 求一阶偏导数，得到得分函数 $S(R_t) \triangleq \frac{\partial \ln L}{\partial R_t}$：
\begin{align}
S(R_t) &= \frac{c}{R_t} - (nk + c) \frac{\frac{\partial}{\partial R_t}(k + \rho R_t)}{k + \rho R_t} \notag \\
&= \frac{c}{R_t} - \frac{\rho(nk + c)}{k + \rho R_t}.
\end{align}
进行通分合并：
\begin{align}
S(R_t) &= \frac{c(k + \rho R_t) - \rho R_t(nk + c)}{R_t(k + \rho R_t)} \notag \\
&= \frac{ck + c \rho R_t - \rho n k R_t - c \rho R_t}{R_t(k + \rho R_t)} \notag \\
&= \frac{k(c - \rho n R_t)}{R_t(k + \rho R_t)}. \label{eq:score_function}
\end{align}
令得分函数等于零 $S(R_t) = 0$，由于 $k > 0$ 且分母在内部空间恒正，驻点方程退化为：
\begin{equation}
c - \rho n R_t = 0 \implies \widehat{R}_t^{\text{MLE}} = \frac{c}{\rho n}.
\end{equation}
又因为二阶导数恒负（见后文 Fisher 信息量推导），该驻点为唯一的全局极大值点，式~\eqref{eq:estimator_def} 证毕。

\textbf{第二步：双层全期望公式（塔式法则）严格无偏性证明}

对估计量 $\widehat{R}_t = \frac{C_t}{\rho n}$ 求条件数学期望：
\begin{equation}
\E\left[ \widehat{R}_t \;\middle|\; n, R_t \right] = \E\left[ \frac{C_t}{\rho n} \;\middle|\; n, R_t \right] = \frac{1}{\rho n} \E\left[ C_t \;\middle|\; n, R_t \right].
\end{equation}
应用条件全期望法则（Law of Total Expectation, 塔式法则），将观测变量 $C_t$ 投影到潜变量 $I_t$ 上：
\begin{equation}
\E\left[ C_t \;\middle|\; n, R_t \right] = \E_{I_t}\Big[ \E_{C_t}\big( C_t \mid I_t \big) \;\Big|\; n, R_t \Big].
\end{equation}
\begin{itemize}
    \item 内层期望：在给定真实感染 $I_t$ 时，抽稀过程服从 $C_t \mid I_t \sim \Binomial(I_t, \rho)$，二项分布均值为 $\E[C_t \mid I_t] = \rho I_t$；
    \item 外层期望：在给定上一代风险集 $n$ 时，真实感染 $I_t = \sum_{j=1}^n X_{t,j}$ 为 $n$ 个均值为 $R_t$ 的独立分支和，故 $\E[I_t \mid n, R_t] = n R_t$。
\end{itemize}
代入得：
\begin{equation}
\E\left[ C_t \;\middle|\; n, R_t \right] = \E_{I_t}\big[ \rho I_t \;\big|\; n, R_t \big] = \rho \E[I_t \mid n, R_t] = \rho(n R_t) = \rho n R_t.
\end{equation}
两端除以常数 $\rho n$：
\begin{equation}
\E\left[ \widehat{R}_t \;\middle|\; n, R_t \right] = \frac{\rho n R_t}{\rho n} \equiv R_t.
\end{equation}
在全参数空间中无偏性严格证毕。

\textbf{第三步：Fisher 信息量双套算法验证与 Cramér--Rao 下界确立}

下面采用数理统计中两种截然不同的定义算法计算 Fisher 信息量 $\mathcal{I}(R_t \mid n)$：

\textit{算法一（得分函数方差法）：}
由得分函数方差定义 $\mathcal{I}(R_t) \triangleq \Var(S(R_t))$。利用式~\eqref{eq:score_function}：
\begin{equation}
\Var\big( S(R_t) \big) = \Var\left( \frac{k(C_t - \rho n R_t)}{R_t(k + \rho R_t)} \right) = \frac{k^2}{R_t^2 (k + \rho R_t)^2} \Var(C_t \mid n, R_t).
\end{equation}
将式~\eqref{eq:obs_var_decomp} 导出的负二项方差 $\Var(C_t \mid n, R_t) = \frac{n \rho R_t(k + \rho R_t)}{k}$ 代入：
\begin{align}
\mathcal{I}(R_t \mid n) &= \frac{k^2}{R_t^2 (k + \rho R_t)^2} \cdot \left[ \frac{n \rho R_t (k + \rho R_t)}{k} \right] \notag \\
&= \frac{k \cdot n \rho R_t (k + \rho R_t)}{R_t^2 (k + \rho R_t)^2} \notag \\
&= \frac{k n \rho}{R_t (k + \rho R_t)}. \label{eq:fisher_alg1}
\end{align}

\textit{算法二（对数似然负二阶导期望法）：}
对得分函数再次求偏导：
\begin{align}
\frac{\partial^2 \ln L}{\partial R_t^2} &= \frac{\partial}{\partial R_t} \left[ \frac{C_t}{R_t} - \frac{\rho(nk + C_t)}{k + \rho R_t} \right] \notag \\
&= -\frac{C_t}{R_t^2} - \rho(nk + C_t) \left[ -\frac{\rho}{(k + \rho R_t)^2} \right] \notag \\
&= -\frac{C_t}{R_t^2} + \frac{\rho^2 (nk + C_t)}{(k + \rho R_t)^2}.
\end{align}
取负号并求数学期望，代入 $\E[C_t \mid n, R_t] = \rho n R_t$：
\begin{align}
\mathcal{I}(R_t \mid n) &= \E\left[ -\frac{\partial^2 \ln L}{\partial R_t^2} \;\middle|\; n, R_t \right] \notag \\
&= \frac{\E[C_t]}{R_t^2} - \frac{\rho^2(nk + \E[C_t])}{(k + \rho R_t)^2} \notag \\
&= \frac{\rho n R_t}{R_t^2} - \frac{\rho^2(nk + \rho n R_t)}{(k + \rho R_t)^2} \notag \\
&= \frac{\rho n}{R_t} - \frac{\rho^2 n(k + \rho R_t)}{(k + \rho R_t)^2} \notag \\
&= \frac{\rho n}{R_t} - \frac{\rho^2 n}{k + \rho R_t} \notag \\
&= \frac{\rho n (k + \rho R_t) - \rho^2 n R_t}{R_t(k + \rho R_t)} \notag \\
&= \frac{k n \rho}{R_t (k + \rho R_t)}. \label{eq:fisher_alg2}
\end{align}
两种独立算法的推导结果在代数上完全恒等！

计算 Fisher 信息量的倒数：
\begin{align}
\mathcal{I}(R_t \mid n)^{-1} &= \frac{R_t(k + \rho R_t)}{k n \rho} \notag \\
&= \frac{k R_t + \rho R_t^2}{k n \rho} \notag \\
&= \frac{k R_t}{k n \rho} + \frac{\rho R_t^2}{k n \rho} \notag \\
&= \frac{R_t}{\rho n} + \frac{R_t^2}{k n}. \label{eq:fisher_inv}
\end{align}
根据经典 Cramér--Rao 定理，对于参数 $R_t$ 的任意无偏估计量 $T(C_t)$，其方差均满足下界：
\begin{equation}
\Var\big( T(C_t) \big) \ge \mathcal{I}(R_t \mid n)^{-1}.
\end{equation}
而对于估计量 $\widehat{R}_t = \frac{C_t}{\rho n}$，直接由方差性质计算：
\begin{equation}
\Var(\widehat{R}_t \mid n, R_t) = \Var\left( \frac{C_t}{\rho n} \right) = \frac{1}{\rho^2 n^2} \Var(C_t \mid n, R_t) = \frac{1}{\rho^2 n^2} \left[ \rho n R_t + \frac{(\rho R_t)^2 n}{k} \right] = \frac{R_t}{\rho n} + \frac{R_t^2}{k n}.
\end{equation}
估计量的方差严格等于 Cramér--Rao 下界，因此 $\widehat{R}_t$ 为 UMVUE，式~\eqref{eq:var_prop1} 严格证毕。
\end{proof}

命题 1 确立了在目标 A（当期瞬态点估计）下的核心结论：**方差随着期望报告规模 $m = \rho n$ 按 $O(1/m)$ 严格无界递减**。在此视角下，大数定律完全起效，监测规模确实“越大越好”。

\section{目标 B 的全方差分解与仿射损失结构}\label{sec:prop2}

本节将推断目标切换至目标 B（跨期平稳基础基准 $\bar{R}$），揭示方差结构的根本突变。

\subsection{不可约宏观环境方差的数理引入}
在现实中，气象变化（如降温导致室内聚集增加）、法定节假日人群大范围迁徙、或者自发佩戴口罩等社会干预，构成全社会所有个体共同面对的外生环境扰动。
设当期瞬态再生数在潜在基础基准 $\bar{R}$ 附近发生宏观乘性扰动：
\begin{equation}
R_t = \bar{R} \cdot \exp(\varepsilon_t), \quad \varepsilon_t \sim \mathcal{N}\left(-\frac{1}{2}v_\varepsilon, \; v_\varepsilon\right),
\end{equation}
满足条件均值 $\E[R_t \mid n] = \bar{R}$，且条件方差由正实数刻画：
\begin{equation}
\Var(R_t \mid n) = v_R > 0.
\end{equation}
这一设定揭示了不可约环境方差 $v_R$ 在两类推断目标下的本质分化：对目标 A 而言，环境波动正是待测对象本身，其跨期变异不构成误差；对目标 B 而言，同一波动退化为不可消除的误差源，单期样本无论多大都无法消除全员共享的单次环境实现。
我们称 $v_R$ 为**不可约环境方差（Irreducible Environmental Variance）**。
\subsection{全方差公式展开与仿射损失结构推导}

\begin{proposition}[\textbf{命题 2：跨期全方差分解与仿射损失}]\label{prop:prop2}
在推断目标 B 下，以单期估计量 $\widehat{R}_t = \frac{C_t}{\rho n}$ 估计跨期基准 $\bar{R}$，在给定风险集 $n$ 条件下，单期边际预测方差严格满足仿射损失结构：
\begin{equation}
\Var\left( \frac{C_t}{\rho n} \;\middle|\; n \right) = v_R + \frac{\bar{R}}{\rho n} + \frac{\bar{R}^2 + v_R}{kn} \equiv v_R + \frac{A}{m}, \label{eq:affine_variance}
\end{equation}
其中期望报告规模为 $m \equiv \rho n$，标量系数 $A$ 由系统深层结构参数唯一定义：
\begin{equation}
A \equiv \bar{R} + \frac{\rho(\bar{R}^2 + v_R)}{k}. \label{eq:scalar_A_def}
\end{equation}
\end{proposition}

\begin{proof}
应用**全方差公式（Law of Total Variance）**：
\begin{equation}
\Var\left( \frac{C_t}{\rho n} \;\middle|\; n \right) = \underbrace{\Var\left( \E\left[ \frac{C_t}{\rho n} \;\middle|\; n, R_t \right] \;\middle|\; n \right)}_{\text{宏观环境波动贡献项}} + \underbrace{\E\left[ \Var\left( \frac{C_t}{\rho n} \;\middle|\; n, R_t \right) \;\middle|\; n \right]}_{\text{微观抽样波动贡献项}}. \label{eq:total_var_law}
\end{equation}

下面分步计算式~\eqref{eq:total_var_law} 中的两项：

\textbf{第一项计算：}
由命题 1 中的无偏性结论，在给定当期环境状态 $R_t$ 时，$\E\left[ \frac{C_t}{\rho n} \;\middle|\; n, R_t \right] = R_t$。代入得：
\begin{equation}
\Var\left( \E\left[ \frac{C_t}{\rho n} \;\middle|\; n, R_t \right] \;\middle|\; n \right) = \Var(R_t \mid n) \equiv v_R.
\end{equation}

\textbf{第二项计算：}
由命题 1 式~\eqref{eq:var_prop1}，已知条件方差为 $\frac{R_t}{\rho n} + \frac{R_t^2}{kn}$。对外层随机变量 $R_t$ 求期望：
\begin{equation}
\E\left[ \frac{R_t}{\rho n} + \frac{R_t^2}{kn} \;\middle|\; n \right] = \frac{\E[R_t \mid n]}{\rho n} + \frac{\E[R_t^2 \mid n]}{kn}.
\end{equation}
利用方差与二阶原点矩的基本代数恒等式：
\begin{equation}
\E[R_t^2 \mid n] = \Var(R_t \mid n) + \big( \E[R_t \mid n] \big)^2 = v_R + \bar{R}^2.
\end{equation}
代入上式即得：
\begin{equation}
\E\left[ \Var\left( \frac{C_t}{\rho n} \;\middle|\; n, R_t \right) \;\middle|\; n \right] = \frac{\bar{R}}{\rho n} + \frac{\bar{R}^2 + v_R}{kn}.
\end{equation}

将两项相加，得到总边际方差：
\begin{equation}
\Var\left( \frac{C_t}{\rho n} \;\middle|\; n \right) = v_R + \frac{\bar{R}}{\rho n} + \frac{\bar{R}^2 + v_R}{kn}.
\end{equation}
代入期望报告规模 $m = \rho n$（即 $n = m / \rho$）：
\begin{align}
\Var\left( \frac{C_t}{\rho n} \;\middle|\; n \right) &= v_R + \frac{\bar{R}}{m} + \frac{\rho(\bar{R}^2 + v_R)}{k m} \notag \\
&= v_R + \frac{1}{m}\left[ \bar{R} + \frac{\rho(\bar{R}^2 + v_R)}{k} \right] \notag \\
&\equiv v_R + \frac{A}{m}.
\end{align}
命题 2 证毕。
\end{proof}

\subsection{特征规模平衡点 $m_\times$ 的内生导出与物理意义}
式~\eqref{eq:affine_variance} 显示，在目标 B 下，边际方差由两部分构成：
\begin{itemize}
    \item 宏观环境方差 $v_R$：与病例样本量 $m$ 完全无关的不可约常数；
    \item 微观抽样方差 $A/m$：随规模 $m$ 按反比例递减的微观离散项。
\end{itemize}
联立这两项方差，自然导出一个两类方差相等的特征交点：
\begin{definition}[\textbf{特征规模平衡点 / 方差平衡点}]\label{def:crossover}
定义抽样方差项与不可约环境方差底板相等的样本量为**特征规模平衡点** $m_\times$：
\begin{equation}
\frac{A}{m_\times} = v_R \implies m_\times \equiv \frac{A}{v_R} = \frac{\bar{R} + \frac{\rho(\bar{R}^2 + v_R)}{k}}{v_R}. \label{eq:crossover_def}
\end{equation}
\noindent\textit{注记：}标量系数 $A \equiv \bar{R} + \frac{\rho(\bar{R}^2 + v_R)}{k}$ 自身仍显式依赖于报告率 $\rho$。因此两个监测系统即使处于相同的 $m = \rho n$，只要报告率 $\rho$ 不同，$A$ 与收益拐点 $m_\times$ 便不相同。$A$ 依赖 $\rho$ 的方式与依赖规模 $m$ 的反比衰减完全不同，这正是后文命题~\ref{prop:prop3} 监测杠杆非对称性的数理根源。
\end{definition}

平衡点 $m_\times$ 具有极强的定性物理分水岭含义：
\begin{itemize}
    \item **当 $m \ll m_\times$ 时（微观噪声主导体制）**：$A/m \gg v_R$，微观抽样与超传播噪声占总方差的 50\% 以上。此时扩大捕获样本量能够带来显著的方差收益；
    \item **当 $m \gg m_\times$ 时（宏观环境饱和体制）**：$A/m \ll v_R$，不可约环境方差占比超过 50\%，总方差曲线进入水平平台。此时继续扩大病例捕获已无法从根本上改善精度，监测资源边际收益剧烈钝化。
\end{itemize}

\section{监测杠杆的理论非对称性定理}\label{sec:prop3}

面对 $v_R + A/m$ 结构，疾控机构有两种潜在的扩量杠杆：其一是投入资源提升报告率 $\rho$（如提高采样率或筛查深度）；其二是扩大监测的地理或人口基数 $n$。本节证明两种杠杆具有本质的非对称性。

\begin{proposition}[\textbf{命题 3：监测杠杆的理论非对称性定理}]\label{prop:prop3}
将单期边际方差表示为关于基数 $n$ 与报告率 $\rho$ 的二元函数：
\begin{equation}
V(n, \rho) \equiv v_R + \frac{\bar{R}}{\rho n} + \frac{\bar{R}^2 + v_R}{kn}.
\end{equation}
该二元函数在极限状态下满足以下非对称收敛性质：
\begin{enumerate}[label=(\roman*)]
    \item **提高报告率极限（$\rho \to 1$）**：
    \begin{equation}
    \lim_{\rho \to 1} V(n, \rho) = v_R + \frac{\bar{R}}{n} + \frac{\bar{R}^2 + v_R}{kn} > v_R. \label{eq:lim_rho}
    \end{equation}
    \item **扩大风险基数极限（$n \to \infty$）**：
    \begin{equation}
    \lim_{n \to \infty} V(n, \rho) = v_R. \label{eq:lim_n}
    \end{equation}
\end{enumerate}
\end{proposition}

\begin{proof}
将第一项抽样项进行代数裂项分解：
\begin{equation}
\frac{\bar{R}}{\rho n} = \frac{\bar{R}}{n} + \frac{(1 - \rho)\bar{R}}{\rho n}. \label{eq:split_decomp}
\end{equation}
则方差函数可精确分解为三层结构：
\begin{equation}
V(n, \rho) = v_R + \underbrace{\frac{(1 - \rho)\bar{R}}{\rho n}}_{\text{漏报抽稀噪声}} + \underbrace{\frac{\bar{R}}{n}}_{\text{有限人群泊松计数波动}} + \underbrace{\frac{\bar{R}^2 + v_R}{kn}}_{\text{超级传播子代抽样波动}}. \label{eq:three_layers}
\end{equation}

\textbf{对式~\eqref{eq:three_layers} 取 $\rho \to 1$ 的极限：}
由于 $\lim_{\rho \to 1} \frac{(1-\rho)\bar{R}}{\rho n} = 0$，漏报抽稀噪声被完全消除。然而，第二项泊松计数波动与第三项超传播波动仅依赖于有限基数 $n$，与 $\rho$ 无关：
\begin{equation}
\lim_{\rho \to 1} V(n, \rho) = v_R + 0 + \frac{\bar{R}}{n} + \frac{\bar{R}^2 + v_R}{kn} = v_R + \frac{1}{n}\left( \bar{R} + \frac{\bar{R}^2 + v_R}{k} \right).
\end{equation}
因为 $\bar{R} > 0, k > 0, n < \infty$，括号内系数严格为正，故极限严格大于 $v_R$。

\textbf{对式~\eqref{eq:three_layers} 取 $n \to \infty$ 的极限（对任意给定的 $\rho > 0$）：}
由于两个规模相关项的分母均含有 $n$：
\begin{equation}
\lim_{n \to \infty} \left[ \frac{\bar{R}}{\rho n} + \frac{\bar{R}^2 + v_R}{kn} \right] = \lim_{n \to \infty} \frac{1}{n} \left[ \frac{\bar{R}}{\rho} + \frac{\bar{R}^2 + v_R}{k} \right] = 0.
\end{equation}
因此：
\begin{equation}
\lim_{n \to \infty} V(n, \rho) = v_R + 0 = v_R.
\end{equation}
命题 3 严格证毕。
\end{proof}

\noindent\textbf{流行病学与管理学启示}：
命题 3 从数理上澄清了一个长期的管理误区：完全消除漏报（$\rho \to 1$）只能消灭由于二项抽样带来的观测抖动，但人群内部由于有限个体相遇产生的泊松随机性以及个体超传播的重尾离散性，是任何高精度检测手段都无法在有限人群中消灭的；唯有扩大监测的地理范围与人群基数 $n \to \infty$，大数定律才能统摄消除全部微观噪声。

\section{多期时间抽样设计与样本预算恒常性}\label{sec:prop4}

当单期横截面样本无法突破环境底板 $v_R$ 时，平滑环境噪声的唯一路径是引入时间维度。本节探讨多期时间抽样设计。

设监测周期跨越 $T$ 期（如连续监测 $T$ 周），在每一期中，疾控系统维持期望报告规模 $m$。跨期样本均值估计量定义为：
\begin{equation}
\overline{\widehat{R}}_T \triangleq \frac{1}{T} \sum_{t=1}^T \widehat{R}_t = \frac{1}{T} \sum_{t=1}^T \frac{C_t}{\rho n_T}.
\end{equation}

\begin{proposition}[\textbf{命题 4：样本预算恒常性与自相关求和}]\label{prop:prop4}
设监测系统在整个评估周期内的总样本人数预算为常数 $B \equiv T \cdot m = T \rho n_T$。
\begin{enumerate}[label=(\roman*)]
    \item 若环境冲击在各期独立，则跨期均值方差为：
    \begin{equation}
    \Var(\overline{\widehat{R}}_T) = \frac{v_R}{T} + \frac{A}{B}. \label{eq:budget_constant_var}
    \end{equation}
    微观抽样误差恒为常数 $A/B$；
    \item 若宏观环境满足一阶自回归协方差 $\Cov(R_t, R_{t+s}) = v_R r^s$（$r \in [0, 1)$），则总方差严格由下式给出：
    \begin{equation}
    \Var(\overline{\widehat{R}}_T) = \frac{v_R}{T}\left[ \frac{1+r}{1-r} - \frac{2r(1-r^T)}{T(1-r)^2} \right] + \frac{A}{B}. \label{eq:ar1_closed_form}
    \end{equation}
    当 $T \gg \frac{1}{1-r}$ 时，有效独立期数渐近退化为 $T_{\mathrm{eff}} \approx T \frac{1-r}{1+r}$。
\end{enumerate}
\end{proposition}

\begin{proof}
应用全方差公式展开：
\begin{equation}
\Var(\overline{\widehat{R}}_T) = \Var\left( \E\left[ \frac{1}{T}\sum_{t=1}^T \widehat{R}_t \;\middle|\; \{R_t\}_{t=1}^T \right] \right) + \E\left[ \Var\left( \frac{1}{T}\sum_{t=1}^T \widehat{R}_t \;\middle|\; \{R_t\}_{t=1}^T \right) \right].
\end{equation}
由单期估计误差在给定环境序列下的条件期望为零及条件独立性：
\begin{equation}
\Var(\overline{\widehat{R}}_T) = \frac{1}{T^2} \Var\left( \sum_{t=1}^T R_t \right) + \frac{1}{T^2} \sum_{t=1}^T \Var(\widehat{R}_t \mid R_t).
\end{equation}
代入单期条件抽样方差 $\Var(\widehat{R}_t \mid R_t) = A/m$：
\begin{equation}
\text{抽样项} = \frac{1}{T^2} \cdot T \left( \frac{A}{m} \right) = \frac{A}{T m}.
\end{equation}
代入总样本预算约束 $B = T m$：
\begin{equation}
\text{抽样项} \equiv \frac{A}{B}.
\end{equation}
此即证明了**样本预算恒常性**：在固定总样本人数配额下，无论如何划分期数 $T$ 与单期规模 $m$，微观抽样误差永远是固定常数 $A/B$！式~\eqref{eq:budget_constant_var} 证毕。

下面给出 AR(1) 环境自相关下双重级数求和的**保姆级错位相减法推导**：

展开环境方差项：
\begin{equation}
\Var\left( \sum_{t=1}^T R_t \right) = \sum_{t=1}^T \Var(R_t) + 2 \sum_{t=1}^{T-1} \sum_{l=t+1}^T \Cov(R_t, R_l).
\end{equation}
代入方差 $v_R$ 与滞后协方差 $v_R r^{l-t}$：
\begin{equation}
\Var\left( \sum_{t=1}^T R_t \right) = T v_R + 2 v_R \sum_{s=1}^{T-1} (T - s) r^s. \label{eq:double_sum}
\end{equation}
定义待求解的等差乘等比级数为：
\begin{equation}
S \triangleq \sum_{s=1}^{T-1} (T - s) r^s = (T-1)r + (T-2)r^2 + (T-3)r^3 + \dots + 1 \cdot r^{T-1}. \label{eq:series_S}
\end{equation}
两端同乘公比 $r$：
\begin{equation}
r S = (T-1)r^2 + (T-2)r^3 + \dots + 2 r^{T-1} + 1 \cdot r^T. \label{eq:series_rS}
\end{equation}
式~\eqref{eq:series_S} 减去式~\eqref{eq:series_rS}，错位相减：
\begin{align}
(1 - r) S &= (T-1)r - \left[ r^2 + r^3 + \dots + r^{T-1} \right] - r^T \notag \\
&= T r - \left[ r + r^2 + r^3 + \dots + r^{T-1} + r^T \right] \notag \\
&= T r - \sum_{s=1}^T r^s.
\end{align}
利用有限等比数列求和公式，$\sum_{s=1}^T r^s = \frac{r(1 - r^T)}{1 - r}$。代入得：
\begin{equation}
(1 - r) S = T r - \frac{r(1 - r^T)}{1 - r}.
\end{equation}
两端除以 $(1 - r)$：
\begin{equation}
S = \frac{T r}{1 - r} - \frac{r(1 - r^T)}{(1 - r)^2}. \label{eq:series_S_solution}
\end{equation}
将式~\eqref{eq:series_S_solution} 回代入总方差式~\eqref{eq:double_sum}：
\begin{align}
\Var\left( \sum_{t=1}^T R_t \right) &= T v_R + 2 v_R \left[ \frac{T r}{1 - r} - \frac{r(1 - r^T)}{(1 - r)^2} \right] \notag \\
&= v_R \left[ T + \frac{2 T r}{1 - r} - \frac{2 r(1 - r^T)}{(1 - r)^2} \right] \notag \\
&= v_R T \left[ 1 + \frac{2r}{1-r} - \frac{2r(1 - r^T)}{T(1-r)^2} \right].
\end{align}
对括号内前两项通分：$1 + \frac{2r}{1-r} = \frac{1-r+2r}{1-r} = \frac{1+r}{1-r}$。因此：
\begin{equation}
\Var\left( \sum_{t=1}^T R_t \right) = v_R T \left[ \frac{1+r}{1-r} - \frac{2r(1 - r^T)}{T(1-r)^2} \right].
\end{equation}
两端除以 $T^2$，与抽样项 $A/B$ 相加，式~\eqref{eq:ar1_closed_form} 严格证毕。
\end{proof}

\section{经费预算约束下的运筹优化与 Cochran 抽样对偶}\label{sec:prop5_cochran}

在现实中，监测期数不可能无限拉长，因为每维持一期系统运转都需要固定的行政与运维经费。本节引入显式经济学成本模型。

设每期维持监测网络的基本固定运维成本为 $c_0 > 0$（如实验室耗材折旧、人员值班底薪），每捕获并检测一例病例的边际成本为 $c_1 > 0$。则在持续 $T$ 期、每期捕获 $m$ 例的设计方案下，总经费消耗为：
\begin{equation}
\mathcal{B} = T (c_0 + c_1 m). \label{eq:cost_budget}
\end{equation}

\begin{corollary}[\textbf{推论 1：Cochran 两阶段抽样对偶与严格凸性}]\label{cor:cochran}
在总经费预算 $\mathcal{B}$ 约束下，假设单期成本在可行域内满足 $\mathcal{B} \ge c_0 + c_1 m$：
\begin{enumerate}[label=(\roman*)]
    \item 跨期目标方差关于单期规模 $m$ 的函数 $f(m) \equiv \Var(\overline{\widehat{R}}_T)$ 在正实轴 $(0, \infty)$ 上严格凸（Strictly Convex）；
    \item 最优单期规模存在唯一的显式解析闭式解：
    \begin{equation}
    m^* = \sqrt{\frac{c_0}{c_1} m_\times} = \sqrt{\frac{c_0 A}{c_1 v_R}}, \label{eq:m_star_closed}
    \end{equation}
    其在数学形式上同构于经典 Cochran (1977) 两阶段抽样最优次级样本量模型；
    \item 若由内部一阶条件求出的规模超出可行上界 $m^* > m_{\max}$，由严格凸性，最优解退化为边界点 $m^* = m_{\max}$。
\end{enumerate}
\end{corollary}

\begin{proof}
由式~\eqref{eq:cost_budget} 反解出期数 $T$：
\begin{equation}
T = \frac{\mathcal{B}}{c_0 + c_1 m}.
\end{equation}
将其代入跨期均值方差公式 $V(m, T) = \frac{v_R}{T} + \frac{A}{T m}$ 中，消去自由变量 $T$：
\begin{align}
f(m) &\triangleq \frac{v_R}{\frac{\mathcal{B}}{c_0 + c_1 m}} + \frac{A}{\frac{\mathcal{B}}{c_0 + c_1 m} \cdot m} \notag \\
&= \frac{c_0 + c_1 m}{\mathcal{B}} \left( v_R + \frac{A}{m} \right) \notag \\
&= \frac{1}{\mathcal{B}} \left( c_0 v_R + c_1 v_R m + \frac{c_0 A}{m} + c_1 A \right). \label{eq:variance_objective}
\end{align}

\textbf{一阶导数与驻点求解：}
对 $m$ 求一阶导数：
\begin{equation}
f'(m) = \frac{1}{\mathcal{B}} \left( c_1 v_R - \frac{c_0 A}{m^2} \right). \label{eq:first_derivative}
\end{equation}
令 $f'(m) = 0$：
\begin{equation}
c_1 v_R - \frac{c_0 A}{m^2} = 0 \implies m^2 = \frac{c_0 A}{c_1 v_R} = \frac{c_0}{c_1} \cdot \left( \frac{A}{v_R} \right).
\end{equation}
代入平衡点定义 $m_\times \equiv A / v_R$，取唯一正实根：
\begin{equation}
m^* = \sqrt{\frac{c_0}{c_1} m_\times}.
\end{equation}

\textbf{二阶导数与严格凸性证明：}
对一阶导数式~\eqref{eq:first_derivative} 再次求导：
\begin{equation}
f''(m) = \frac{1}{\mathcal{B}} \frac{d}{dm}\left( c_1 v_R - c_0 A m^{-2} \right) = \frac{1}{\mathcal{B}} \left( - c_0 A \cdot (-2 m^{-3}) \right) = \frac{2 c_0 A}{\mathcal{B} m^3}. \label{eq:second_derivative}
\end{equation}
因为固定成本 $c_0 > 0$、总经费 $\mathcal{B} > 0$、方差标量 $A > 0$，因此当规模 $m \in (0, \infty)$ 时：
\begin{equation}
f''(m) > 0 \quad (\forall m > 0).
\end{equation}
二阶导数恒严格为正，数学上严格证明了目标函数 $f(m)$ 在正半轴上为**严格凸函数**。因此，驻点 $m^*$ 必然是全局唯一的最小值点！

\textbf{Cochran 抽样理论同构：}
在经典调查抽样理论中，Cochran (1977) 第 10 章给出的两阶段抽样模型为：在群成本 $c_1$ 与人均成本 $c_2$ 下，最小化估计方差的每群次级样本量公式为 $m_{\text{Cochran}} = \sqrt{\frac{c_1 S_2^2}{c_2 (S_1^2 - S_2^2/M)}}$。
在本模型中：
\begin{itemize}
    \item 初级抽样单元（PSU）精确对应时间跨度期数 $T$，单期固定维持成本对应群成本 $c_0$；
    \item 次级抽样单元（SSU）精确对应单期病例样本规模 $m$，单例检测成本对应边际成本 $c_1$；
    \item 群间方差对应宏观环境方差 $v_R$，群内方差对应微观超传播方差 $A$。
\end{itemize}
二者在代数结构上在代数形式上完全同构。推论 1 证毕。
\end{proof}
